% BHU PYQ -- Manually transcribed & error-corrected
% Paper: BCH-225 : Public Finance
% Exam: B.Com. (Hons.) Semester IV Examination, 2022-23

\documentclass[11pt,a4paper]{article}
\usepackage[a4paper,top=2.5cm,bottom=2.5cm,left=2.8cm,right=2.8cm,headheight=14pt]{geometry}
\usepackage{amsmath,amssymb}
\usepackage[T1]{fontenc}
\usepackage[utf8]{inputenc}
\usepackage{lmodern,microtype,fancyhdr,enumitem,hyperref,lastpage,parskip}

\hypersetup{
    colorlinks=true,
    urlcolor=black,
    linkcolor=black,
    pdftitle={BCH-225: Public Finance -- BHU B.Com. (Hons.) Sem IV 2022-23}
}

\pagestyle{fancy}
\fancyhf{}
\renewcommand{\headrulewidth}{0.4pt}
\renewcommand{\footrulewidth}{0.4pt}
\fancyhead[L]{\small\textit{Banaras Hindu University}}
\fancyhead[R]{\small\textit{BCH-225: Public Finance}}
\fancyfoot[C]{\small Page \thepage\ of \pageref{LastPage}}

\newlist{parts}{enumerate}{2}
\setlist[parts,1]{label=(\alph*),leftmargin=*,itemsep=5pt,topsep=3pt}
\setlist[parts,2]{label=(\roman*),leftmargin=*,itemsep=3pt,topsep=2pt}
\newcommand{\pts}[1]{\hfill\textit{[#1]}}

\begin{document}

\begin{center}
  {\large\bfseries BANARAS HINDU UNIVERSITY}\\[2pt]
  {\normalsize B.Com. (Hons.) --- Semester IV Examination, 2022-23}\\[6pt]
  \rule{0.8\linewidth}{0.5pt}\\[6pt]
  {\large\bfseries Paper: BCH-225: Public Finance}\\[4pt]
  {\normalsize Time: 3 Hours\hfill Maximum Marks: 70}
\end{center}

\rule{\linewidth}{0.4pt}

\smallskip
\noindent\textit{Instructions: Answer all questions. All questions carry equal marks (14 marks each).}

\bigskip

\noindent\textbf{Question 1.} \\
Define Public Finance. How does it relate to other sciences? Discuss.\pts{14}

\centerline{\textbf{OR}}

What are the similarities and dis-similarities between public finance and private finance? Explain.\pts{14}

\medskip\rule{\linewidth}{0.2pt}

\noindent\textbf{Question 2.} \\
What is public expenditure? Examine the effects of public expenditure on production and distribution.\pts{14}

\centerline{\textbf{OR}}

How does public expenditure differ from private expenditure? Discuss.\pts{14}

\medskip\rule{\linewidth}{0.2pt}

\noindent\textbf{Question 3.} \\
Discuss the subjective and objective approaches to ``Ability to Pay Tax Principle'' and show how these are applicable in practice.\pts{14}

\centerline{\textbf{OR}}

Discuss Adam Smith's canons of taxation. What additions have been made to them by later economists?\pts{14}

\medskip\rule{\linewidth}{0.2pt}

\noindent\textbf{Question 4.} \\
What is incidence of Tax? Discuss the shifting of incidence of a tax under monopoly market conditions.\pts{14}

\centerline{\textbf{OR}}

What is relative taxable capacity? Discuss the factors affecting the taxable capacity of a country.\pts{14}

\medskip\rule{\linewidth}{0.2pt}

\noindent\textbf{Question 5.} \\
What do you mean by Public Debt? For what purposes can Public Debt be taken? Distinguish between private debt and public debt.\pts{14}

\centerline{\textbf{OR}}

Write short notes on any two of the following:\pts{7+7}
\begin{parts}
  \item Methods of Deficit Financing
  \item Tax Capitalization
  \item Sources of Public Debts
  \item Progressive Tax System
\end{parts}

\vfill
\rule{\linewidth}{0.4pt}
\begin{center}\small
  \href{https://bhu-exam.vercel.app/}{Click here to visit our website}\\[4pt]
  \href{https://me-aryan.vercel.app/}{Click here to visit our contributor}\\[4pt]
  \href{https://quashan-me.vercel.app/}{Click here to visit our contributor}
\end{center}

\end{document}
